\section{Conclusion}
\label{sec:conclusion}

Modular-robot simulation requires a coherent account of meaning as well as mechanics.
A physics engine can integrate bodies and solve a weld, but it does not inherently know
which faces are compatible, which orientation was selected, whether policy permits a
dock, when the logical graph should change, or how an interface should explain the
event. \modsim{} supplies that missing semantic layer while deliberately delegating
physics to backend adapters.

The implemented architecture connects a strict portable \robotpack{}, schema-aware
Studio authoring, canonical runtime state, structured docking acceptance, two-phase
physical/logical commit, deterministic events, derived assemblies, atomic joint
commands, backend capability checks, immutable model views, and a dual-window Runtime
Inspector. A dependency-free mock provides executable semantic reference behavior, and
the \mujoco{} adapter adds articulated dynamics, effort actuation, measured connector
sites, contact, runtime welds, paired collision exclusions, and a native viewer.

The \smoresep{} case study demonstrates the value of keeping these layers separate.
The same connector and topology semantics support a fast kinematic lifecycle, a
two-module physical differential-drive dock/undock cycle, and a seven-module
wheel/contact realization of the published Driver-to-Snake connector plan. The physical
reference run performs four $6\rightarrow5\rightarrow6$ topology transitions, reaches
the exact final chain, records ten docks and four undocks without semantic failure, and
does so without root motion injection after initialization. A fresh current-head
verification run passes 537 tests with 86\% branch-aware core coverage.

The strongest conclusion is architectural, not a claim of hardware fidelity: modular
semantics, physical constraints, canonical events, and generated views can remain
coherent while physical fidelity is increased incrementally. The remaining work is
substantial---calibration, magnetic capture, force reporting, autonomous planning,
additional views, richer Studio/runtime interfaces, a second dynamic backend, and
hardware validation---but it can proceed without replacing the foundation established
here.

